% ============================================================
% ParetoComm: Springer LNNS Conference Paper
% Compile: pdflatex -> bibtex -> pdflatex x2
% ============================================================
\documentclass[runningheads]{llncs}

\usepackage{graphicx}
\usepackage{amsmath,amssymb,amsfonts}
\usepackage{algorithm}
\usepackage{algpseudocode}
\usepackage{booktabs}
\usepackage{multirow}
\usepackage{xcolor}
\usepackage{hyperref}
\usepackage{cite}
\usepackage{bm}
\usepackage{bbm}

\hypersetup{colorlinks=true, linkcolor=blue, citecolor=blue, urlcolor=blue}

% Macros
\newcommand{\E}{\mathbb{E}}
\newcommand{\R}{\mathbb{R}}
\newcommand{\PF}{\mathcal{PF}}
\newcommand{\W}{\mathcal{W}}
\newcommand{\sg}{\mathrm{sg}}
\newcommand{\bw}{\mathbf{w}}
\newcommand{\bG}{\mathbf{G}}
\newcommand{\br}{\mathbf{r}}

\begin{document}

\title{ParetoComm: Learning Sparse Communication Policies\\
on the Multi-Objective Frontier for\\
Resource-Constrained Multi-Agent Systems}

\author{Anandsagar\inst{1} \and
Abhinav Pratap Rai\inst{1}}

\authorrunning{Anandsagar and Rai}

\institute{Department of Computer Science and Engineering,\\
Rajiv Gandhi Institute of Petroleum Technology, Amethi, India\\
\email{\{cs\_student, abhinav.rai\}@rgipt.ac.in}}

\maketitle

% ============================================================
\begin{abstract}
Communication in multi-agent reinforcement learning (MARL) systems
enables coordination but incurs non-trivial resource costs in
bandwidth and energy, especially in edge-deployed systems such as
drone swarms, mobile robots, and wireless sensor networks.
Existing communication-efficient MARL approaches reduce communication
cost by incorporating it as a scalar penalty, which forces a fixed
trade-off and cannot reveal the full spectrum of optimal policies.
We present \textbf{ParetoComm}, a framework that treats communication
cost as an explicit, separate objective alongside task performance,
enabling agents to learn a set of Pareto-optimal communication
policies spanning different coordination--efficiency trade-offs.
ParetoComm extends the Decentralized Partially Observable MDP
(Dec-POMDP) formalism with a binary communication gating action and
a 2-dimensional vector reward.
A learned gating network, trained using the Straight-Through
Gumbel-Sigmoid estimator, enables differentiable binary gating.
We approximate the Pareto front via linear scalarization with a
uniform weight grid, training one policy per weight vector using an
actor-critic update with a shared value baseline.
Experiments on the Multi-Robot Warehouse (RWARE) benchmark with
4 agents demonstrate that ParetoComm achieves a Hypervolume of
$0.71$ ($\pm 0.04$), outperforming a scalar-penalized MAPPO baseline
($0.53 \pm 0.06$) and revealing trade-off regions that scalar
approaches systematically miss.
\keywords{Multi-agent reinforcement learning \and
Multi-objective optimization \and Communication efficiency \and
Pareto front \and Sparse communication}
\end{abstract}

% ============================================================
\section{Introduction}
\label{sec:intro}

Cooperative multi-agent reinforcement learning (MARL) has achieved
remarkable results in complex coordination tasks including autonomous
vehicle platooning~\cite{zhou2019smarts}, cooperative robotic
manipulation~\cite{guhur2021humanoid}, and distributed
sensing~\cite{zhang2021multi}.
A common approach is to augment individual agent policies with
inter-agent communication, allowing agents to share local
observations and improve global coordination~\cite{sukhbaatar2016learning,
das2019tarmac}.

However, in resource-constrained real-world deployments---drone swarms
operating on battery power, IoT sensor networks with limited bandwidth,
and edge-deployed robots with narrow communication channels---every
transmitted message incurs a measurable cost.
Excessive communication leads to energy depletion, network congestion,
and increased latency~\cite{lin2021learning}.
Critically, this cost cannot be ignored or absorbed into the task reward
without losing information about where on the coordination--efficiency
trade-off spectrum a given policy sits.

\paragraph{The limitation of scalar penalization.}
Existing sparse communication methods---such as
CASEC~\cite{wang2020learning}, SchedNet~\cite{kim2019schednet}, and
CommNet with bandwidth constraints~\cite{sukhbaatar2016learning}---
handle communication cost by adding a penalty term
$r_\text{scalar} = r_\text{task} - \lambda \cdot c_\text{comm}$,
where $\lambda$ is a fixed trade-off coefficient.
This reduces a fundamentally multi-objective problem to a single
objective.
The consequence is that training one policy for a fixed $\lambda$
collapses the trade-off to a single point: a practitioner who needs a
different balance must retrain entirely, and has no systematic way to
compare policies across the trade-off spectrum.

\paragraph{Our approach.}
We propose \textbf{ParetoComm}, a framework that explicitly models the
coordination--communication trade-off as a \emph{multi-objective}
problem.
Each agent maintains a binary communication gate, trained with the
Straight-Through Gumbel-Sigmoid estimator~\cite{jang2017categorical}
to produce differentiable discrete gating decisions.
We model the environment as a Dec-POMDP with a 2-dimensional vector
reward and approximate the Pareto front of optimal policies via linear
scalarization with a uniform weight grid.

\paragraph{Contributions.}
\begin{enumerate}
  \item A formal extension of the Dec-POMDP tuple with a decomposed
    action space $A^i = A^i_\text{env} \times A^i_\text{comm}$ and
    a 2-dimensional vector reward.
  \item A learned binary gating mechanism using the Straight-Through
    Gumbel-Sigmoid estimator, ensuring consistent gating behavior
    between training and execution.
  \item A Pareto front approximation framework for cooperative MARL,
    producing a set of policies that outperforms scalar-penalized
    baselines in Hypervolume on the RWARE benchmark.
  \item A complete, open-source implementation enabling reproducible
    evaluation of communication--efficiency trade-offs.
\end{enumerate}

The remainder of this paper is organized as follows.
Section~\ref{sec:related} reviews related work.
Section~\ref{sec:formulation} presents the formal problem setting.
Section~\ref{sec:method} describes the ParetoComm methodology.
Section~\ref{sec:experiments} details the experimental setup.
Section~\ref{sec:results} presents results.
Section~\ref{sec:discussion} discusses limitations.
Section~\ref{sec:conclusion} concludes.

% ============================================================
\section{Related Work}
\label{sec:related}

\subsection{Communication in Multi-Agent RL}

CommNet~\cite{sukhbaatar2016learning} introduced continuous
communication channels where agents broadcast differentiable message
vectors, enabling cooperative behavior but with no mechanism to
suppress unnecessary messages.
TarMAC~\cite{das2019tarmac} added targeted communication via
attention-based routing, improving efficiency but still treating
communication as a soft, continuous process rather than a discrete
resource decision.
ATOC~\cite{jiang2018learning} proposed a binary attention unit that
gates communication based on an agent's internal confidence estimate,
which is closest in spirit to our gating mechanism.
However, all of the above optimize a single scalar objective.
SchedNet~\cite{kim2019schednet} and CASEC~\cite{wang2020learning}
introduced scheduling and sparsity constraints, respectively, but both
incorporate communication cost as a scalar penalty---precluding
systematic trade-off analysis.

\subsection{Multi-Objective Reinforcement Learning}

Multi-Objective RL (MORL) replaces a scalar reward with a vector
$\br_t \in \R^m$ and aims to approximate the Pareto
front~\cite{hayes2022practical}.
Linear scalarization~\cite{roijers2013survey}---training one policy
per weight vector $\bw$---is the simplest approach but is known to
recover only the convex hull of the Pareto
front~\cite{das1997nonlinear}.
Envelope Q-learning~\cite{yang2019generalized} and Pareto
Q-learning~\cite{van2014multi} address non-convex fronts but require
significant architectural changes.
Almost all MORL methods are designed for single-agent settings.
Zintgraf et al.~\cite{zintgraf2015quality} extended MORL to
multi-agent coordination but did not address learned communication.

\subsection{Research Gap}

To the best of our knowledge, no existing work combines learned binary
communication gating with explicit multi-objective optimization in a
decentralized cooperative MARL setting.
ParetoComm bridges this gap.

% ============================================================
\section{Problem Formulation}
\label{sec:formulation}

\subsection{Dec-POMDP with Communication Extension}

We model the environment as a \textbf{Decentralized Partially Observable
MDP (Dec-POMDP)}~\cite{bernstein2002complexity} with the following tuple:
\begin{equation}
  \langle N, S, \{A^i\}, \{O^i\}, T, R, \gamma, \mu_0 \rangle
  \label{eq:decpomdp}
\end{equation}
where:
\begin{itemize}
  \item $N = \{1, \ldots, n\}$ is the set of agents.
  \item $S$ is the global state space (unobservable by any single agent).
  \item $A^i$ is the action space of agent $i$ (defined below).
  \item $O^i$ is the local observation space of agent $i$; agent $i$
    observes $o^i_t \sim \mathcal{O}(s_t)$ but cannot directly access $s_t$.
  \item $T: S \times \prod_i A^i \to \Delta(S)$ is the transition kernel.
  \item $R: S \times \prod_i A^i \to \R^2$ is the \emph{vector} reward
    function (Section~\ref{sec:reward}).
  \item $\gamma \in [0,1)$ is the discount factor.
  \item $\mu_0 \in \Delta(S)$ is the initial state distribution.
\end{itemize}

\paragraph{Action space decomposition.}
The key extension is decomposing $A^i$ into two independent components:
\begin{equation}
  A^i = A^i_\text{env} \times A^i_\text{comm},
  \quad A^i_\text{comm} \in \{0, 1\}
  \label{eq:action_decomp}
\end{equation}
At each timestep $t$, agent $i$ selects a tuple
$(a^i_{\text{env},t},\; a^i_{\text{comm},t}) \in A^i_\text{env} \times \{0,1\}$.
The environment transition depends \emph{only} on $a^i_\text{env}$;
the communication action $a^i_\text{comm}$ determines whether a message
vector is transmitted to other agents.
When $a^i_{\text{comm},t} = 0$, agent $i$ broadcasts a zero vector and
incurs no communication cost.

\subsection{Vector Reward Function}
\label{sec:reward}

We use a \textbf{2-dimensional vector reward}:
\begin{equation}
  \br^i_t =
  \begin{bmatrix} r^{\text{task}}_t \\ r^{\text{comm}}_t \end{bmatrix}
  \in \R^2
  \label{eq:reward_vector}
\end{equation}
In our cooperative setting, all agents share the same reward vector.
The components are:

\textbf{Task reward} $r^{\text{task}}_t$ is environment-defined.
In RWARE, it is the number of shelves successfully delivered to the
workstation at timestep $t$: $r^{\text{task}}_t \in \{0, 1\}$.

\textbf{Communication reward} $r^{\text{comm}}_t$ penalizes the total
number of messages sent:
\begin{equation}
  r^{\text{comm}}_t = -\sum_{i=1}^{n} a^i_{\text{comm},t}
  \label{eq:comm_reward}
\end{equation}
This is well-defined for discrete gating decisions: every transmitted
message costs exactly $1$ unit, and no message costs $0$.
The range is $r^{\text{comm}}_t \in \{-n, \ldots, 0\}$.

\paragraph{Why not add an energy term?}
An energy term $r^{\text{energy}}_t \propto \|a^i_{\text{env}}\|$ is
ill-defined when $a^i_\text{env}$ is a discrete action index, as it is
in RWARE.
We therefore use a 2-objective formulation; extending to 3 objectives
with a continuous-action environment is left to future work.

\subsection{Discounted Vector Return}

The \textbf{vector return} over an episode of length $T$ is:
\begin{equation}
  \bG = \sum_{t=0}^{T} \gamma^t \, \br_t \in \R^2
  \label{eq:vector_return}
\end{equation}
This is the object we optimize: $\bG = [G^\text{task},\; G^\text{comm}]^\top$.

\subsection{Pareto Dominance and Pareto Front}

Policy $\pi$ \textbf{dominates} policy $\pi'$ (written $\pi \succ \pi'$)
if and only if:
\begin{equation}
  \forall k \in \{1,2\}:\; \E[G^k_\pi] \geq \E[G^k_{\pi'}]
  \quad \text{and} \quad
  \exists k:\; \E[G^k_\pi] > \E[G^k_{\pi'}]
  \label{eq:dominance}
\end{equation}
The \textbf{Pareto front} is the set of all non-dominated policies:
\begin{equation}
  \PF = \{\pi \mid \nexists\, \pi' \text{ such that } \pi' \succ \pi\}
  \label{eq:pareto_front}
\end{equation}
Our objective is to approximate $\PF$ with a finite set of trained
policies.

\subsection{Scalarized Optimization Objective}

We approximate $\PF$ via linear scalarization.
For a weight vector $\bw = (w_1, w_2) \in \Delta^1$ where
$\Delta^1 = \{\bw \in \R^2 : w_1 + w_2 = 1,\; w_k \geq 0\}$,
the scalarized objective is:
\begin{equation}
  \pi^*_{\bw} = \arg\max_\pi\; \bw^\top \E[\bG_\pi]
  = \arg\max_\pi\; w_1 \E[G^\text{task}_\pi] + w_2 \E[G^\text{comm}_\pi]
  \label{eq:scalarized}
\end{equation}
We train one policy $\pi_{\bw}$ per weight vector $\bw \in \W$, where:
\begin{equation}
  \W = \{(w, 1-w) : w \in \{0.0, 0.1, 0.2, \ldots, 1.0\}\}
  \label{eq:weight_grid}
\end{equation}
giving 11 weight vectors and 11 trained policies.
The Pareto front approximation is:
\begin{equation}
  \hat{\PF} = \{\E[\bG_{\pi_{\bw}}] : \bw \in \W\}
  \label{eq:pf_approx}
\end{equation}

\textbf{Limitation:} Linear scalarization with positive weights
recovers only the \emph{convex hull} of $\PF$~\cite{das1997nonlinear}.
Non-convex regions of the true Pareto front are not guaranteed to be
found. Chebyshev scalarization~\cite{miettinen1999nonlinear} or
envelope Q-learning~\cite{yang2019generalized} address this limitation.

% ============================================================
\section{Methodology}
\label{sec:method}

\subsection{System Overview}

ParetoComm trains a separate policy $\pi_{\bw}$ for each weight
vector $\bw \in \W$.
Each policy consists of shared-parameter agents, where each agent
runs: (1) an observation encoder, (2) a communication gating network
that outputs a binary gate, (3) a message encoder, (4) a message
aggregator, and (5) an action policy with value head.
Figure~\ref{fig:arch} illustrates the per-agent architecture.

\begin{figure}[t]
  \centering
  \includegraphics[width=0.85\textwidth]{figures/architecture.pdf}
  \caption{Per-agent architecture in ParetoComm.
    Dashed connections indicate paths that are suppressed when
    $g^i_t = 0$ (gate closed). The value head is only used during
    training.}
  \label{fig:arch}
\end{figure}

\subsection{Agent Architecture}

Each agent $i$ consists of the following components.

\textbf{Observation encoder} $f_\phi: O^i \to \R^{d_h}$ is a
two-layer MLP that maps the local observation $o^i_t \in \R^{d_o}$
to a hidden representation $h^i_t = f_\phi(o^i_t) \in \R^{d_h}$.

\textbf{Gating network} $f_\theta: \R^{d_h} \to [0,1]$ maps $h^i_t$
to a gate probability $p^i_t = \sigma(f_\theta(h^i_t))$,
where $\sigma$ is the sigmoid function.

\textbf{Message encoder} $\text{Enc}_\psi: \R^{d_h} \to \R^{d_m}$
transforms the hidden representation into a message vector.
The gated message is:
\begin{equation}
  m^i_t = g^i_t \cdot \text{Enc}_\psi(h^i_t)
  \label{eq:gated_message}
\end{equation}
When $g^i_t = 0$, $m^i_t = \mathbf{0}$ (zero vector; no information
transmitted and no cost incurred).

\textbf{Message aggregator} produces agent $i$'s received context by
mean-pooling messages from all other agents:
\begin{equation}
  \bar{m}^i_t = \frac{1}{n-1} \sum_{j \neq i} m^j_t
  \label{eq:aggregator}
\end{equation}

\textbf{Action policy network} $\pi_\omega: \R^{d_h} \times \R^{d_m} \to \Delta(A^i_\text{env})$
takes the concatenation $[h^i_t;\, \bar{m}^i_t]$ and outputs a
distribution over environment actions.

\textbf{Value network} $V_\nu: \R^{d_h} \times \R^{d_m} \to \R$
estimates the scalar value of the current observation--message pair
under scalarized reward $r^\text{scalar}_t = \bw^\top \br_t$.
This is used only during training for variance reduction.

\subsection{Communication Gating with Straight-Through Gumbel-Sigmoid}
\label{sec:gating}

\paragraph{Gate probability.}
\begin{equation}
  p^i_t = \sigma(f_\theta(h^i_t)) \in (0,1)
  \label{eq:gate_prob}
\end{equation}

\paragraph{Hard gate at execution.}
At test time, agent $i$ samples a hard binary gate:
\begin{equation}
  g^i_t \sim \text{Bernoulli}(p^i_t) \in \{0,1\}
  \label{eq:hard_gate}
\end{equation}

\paragraph{Differentiable relaxation at training.}
Since the Bernoulli sample is non-differentiable, we use the
\textbf{Gumbel-Sigmoid} (binary Gumbel-Softmax)~\cite{jang2017categorical,
maddison2017concrete} relaxation during training:
\begin{equation}
  \tilde{g}^i_t = \sigma\!\left(\frac{\log p^i_t - \log(1-p^i_t) + \epsilon_1 - \epsilon_2}{\tau}\right)
  \label{eq:gumbel_sigmoid}
\end{equation}
where $\epsilon_1, \epsilon_2 \sim \text{Gumbel}(0,1)$ are independent
noise samples and $\tau > 0$ is the temperature.
As $\tau \to 0$, $\tilde{g}^i_t \to \text{Bernoulli}(p^i_t)$.

\paragraph{Straight-Through Estimator.}
To ensure the policy uses \emph{hard} binary gates during both training
forward passes and execution, while allowing gradients to flow through
the soft gate during the backward pass, we use the
\textbf{Straight-Through Gumbel estimator}~\cite{jang2017categorical}:
\begin{equation}
  g^i_t = g^i_\text{hard} + \sg(\tilde{g}^i_t - g^i_\text{hard})
  \label{eq:st_estimator}
\end{equation}
where $g^i_\text{hard} = \mathbbm{1}[\tilde{g}^i_t > 0.5] \in \{0,1\}$
and $\sg(\cdot)$ denotes the stop-gradient operator.
In the forward pass, $g^i_t = g^i_\text{hard}$ (hard binary decision).
In the backward pass, gradients flow through $\tilde{g}^i_t$, which is
differentiable with respect to $\theta$.

\paragraph{Temperature annealing.}
We anneal $\tau$ over the course of training using an exponential
schedule based on cumulative environment timesteps $t$:
\begin{equation}
  \tau_t = \max(\tau_\text{min},\; \tau_0 \cdot \lambda^{\lfloor t / \Delta \rfloor})
  \label{eq:annealing}
\end{equation}
where $\tau_0 = 1.0$, $\lambda = 0.995$, $\tau_\text{min} = 0.1$,
and $\Delta = 1000$ (decay applied every 1000 timesteps).
This schedule ensures the gate gradually transitions from soft
(exploratory) to hard (decisive) as training progresses.

\subsection{Actor-Critic Training}

For a given weight vector $\bw$, we define the scalar reward as:
\begin{equation}
  r^\text{scalar}_t = \bw^\top \br_t = w_1 r^\text{task}_t + w_2 r^\text{comm}_t
  \label{eq:scalar_reward}
\end{equation}
and the scalar return as $G^\text{scalar} = \sum_{t=0}^T \gamma^t r^\text{scalar}_t$.

We train using an actor-critic update with a shared value function
$V_\nu$ to reduce gradient variance.
The advantage estimate is:
\begin{equation}
  A_t = G^\text{scalar}_t - V_\nu(o^i_t, \bar{m}^i_t)
  \label{eq:advantage}
\end{equation}
where $G^\text{scalar}_t = \sum_{t'=t}^T \gamma^{t'-t} r^\text{scalar}_{t'}$
is the Monte Carlo return from step $t$.

The \textbf{policy loss} is:
\begin{equation}
  \mathcal{L}_\text{policy}(\omega, \theta) =
  -\E\!\left[\sum_t A_t \left(\log \pi_\omega(a^i_t \mid o^i_t, \bar{m}^i_t)
    + \log p(g^i_t \mid p^i_t)\right)\right]
  \label{eq:policy_loss}
\end{equation}
The \textbf{value loss} is:
\begin{equation}
  \mathcal{L}_\text{value}(\nu) =
  \E\!\left[\sum_t \left(G^\text{scalar}_t - V_\nu(o^i_t, \bar{m}^i_t)\right)^2\right]
  \label{eq:value_loss}
\end{equation}
The \textbf{entropy bonus} encourages exploration:
\begin{equation}
  \mathcal{L}_\text{entropy} =
  -\beta_e \E\!\left[\mathcal{H}(\pi_\omega(\cdot \mid o^i_t, \bar{m}^i_t))\right]
  - \beta_g \E\!\left[\mathcal{H}(p^i_t)\right]
  \label{eq:entropy_loss}
\end{equation}
where $\mathcal{H}(p) = -p\log p - (1-p)\log(1-p)$ for the gate entropy.

The combined loss is:
\begin{equation}
  \mathcal{L} = \mathcal{L}_\text{policy} + c_v \mathcal{L}_\text{value}
    + \mathcal{L}_\text{entropy}
  \label{eq:total_loss}
\end{equation}
with $c_v = 0.5$.

\subsection{Algorithm}

Algorithm~\ref{alg:paretocomm} presents the full training procedure.

\begin{algorithm}[t]
\caption{ParetoComm Training}
\label{alg:paretocomm}
\begin{algorithmic}[1]
\Require Weight grid $\W$, episodes $N$, discount $\gamma$,
         annealing parameters $\tau_0, \lambda, \tau_\text{min}, \Delta$
\Ensure Pareto front approximation $\hat{\PF}$

\State Initialize $\hat{\PF} \leftarrow \emptyset$
\For{each $\bw = (w_1, w_2) \in \W$}
  \State Initialize policy $\pi_\omega$, gate $f_\theta$,
         value $V_\nu$, encoder $f_\phi$, msg encoder $\text{Enc}_\psi$
  \State $t_\text{global} \leftarrow 0$
  \For{episode $= 1$ to $N$}
    \State $\{o^i_0\} \leftarrow \text{env.reset()}$
    \State Initialize trajectory buffer $\mathcal{B} \leftarrow \emptyset$
    \For{$t = 0, 1, \ldots, T-1$}
      \State $\tau \leftarrow \max(\tau_\text{min},\;\tau_0 \cdot \lambda^{\lfloor t_\text{global}/\Delta\rfloor})$
      \For{each agent $i \in N$}
        \State $h^i_t \leftarrow f_\phi(o^i_t)$
        \Comment{Encode observation}
        \State $p^i_t \leftarrow \sigma(f_\theta(h^i_t))$
        \Comment{Gate probability}
        \State $\tilde{g}^i_t \leftarrow \text{GumbelSigmoid}(p^i_t, \tau)$
        \Comment{Eq.~\ref{eq:gumbel_sigmoid}}
        \State $g^i_t \leftarrow g^i_\text{hard} + \sg(\tilde{g}^i_t - g^i_\text{hard})$
        \Comment{Straight-Through, Eq.~\ref{eq:st_estimator}}
        \State $m^i_t \leftarrow g^i_t \cdot \text{Enc}_\psi(h^i_t)$
        \Comment{Gated message, Eq.~\ref{eq:gated_message}}
      \EndFor
      \For{each agent $i \in N$}
        \State $\bar{m}^i_t \leftarrow \frac{1}{n-1}\sum_{j \neq i} m^j_t$
        \Comment{Aggregate messages, Eq.~\ref{eq:aggregator}}
        \State $a^i_t \sim \pi_\omega(\cdot \mid h^i_t, \bar{m}^i_t)$
        \Comment{Sample action}
      \EndFor
      \State $\{o^i_{t+1}\},\; \br_t \leftarrow \text{env.step}(\{a^i_t\})$
      \State $r^\text{scalar}_t \leftarrow \bw^\top \br_t$
      \Comment{Scalarize, Eq.~\ref{eq:scalar_reward}}
      \State $\mathcal{B} \leftarrow \mathcal{B} \cup (o_t, g_t, a_t, r^\text{scalar}_t, o_{t+1})$
      \State $t_\text{global} \leftarrow t_\text{global} + 1$
    \EndFor
    \State Compute returns $\{G^\text{scalar}_t\}$ and advantages $\{A_t\}$ from $\mathcal{B}$
    \Comment{Eq.~\ref{eq:advantage}}
    \State Update $(\omega, \theta, \nu, \phi, \psi)$ by minimizing $\mathcal{L}$
    \Comment{Eq.~\ref{eq:total_loss}}
  \EndFor
  \State Evaluate $\pi_\omega$ over $M$ episodes; record $\E[\bG_{\pi_\bw}]$
  \State $\hat{\PF} \leftarrow \hat{\PF} \cup \{\E[\bG_{\pi_\bw}]\}$
\EndFor
\State Remove dominated points from $\hat{\PF}$
\Comment{Eq.~\ref{eq:dominance}}
\State \Return $\hat{\PF}$
\end{algorithmic}
\end{algorithm}

% ============================================================
\section{Experimental Setup}
\label{sec:experiments}

\subsection{Environment}

We evaluate on the \textbf{Multi-Robot Warehouse (RWARE)}
benchmark~\cite{christianos2020shared}, specifically the
\texttt{rware-tiny-4ag-v1} variant (4 agents, $10 \times 11$ grid,
2 shelves to deliver per episode).
RWARE provides a challenging partially observable cooperative task
where agents must navigate, pick up shelves, and deliver them to a
workstation, with realistic inter-agent collision dynamics.

\textbf{Observation space:} Each agent observes a $\{0,1\}^{d_o}$
binary vector encoding: agent's position and direction, nearby agents,
shelf locations within sensor range, and goal positions.
$d_o = 13$ for the tiny variant.

\textbf{Action space:} $A^i_\text{env} = \{\text{NOOP, FORWARD, LEFT, RIGHT, TOGGLE\_LOAD}\}$,
$|A^i_\text{env}| = 5$.

\textbf{Task reward:} $r^\text{task}_t = $ number of successful deliveries
at timestep $t$; $\in \{0, 1, \ldots, n\}$.

\subsection{Baselines}

\textbf{B1 — Always communicate (MAPPO):} Standard MAPPO~\cite{yu2021surprising}
with $a^i_\text{comm} \equiv 1$ and task reward only.
No communication cost modeled.

\textbf{B2 — Scalar penalized (SP-MARL):} MAPPO with scalar reward
$r^\text{scalar}_t = r^\text{task}_t - \lambda r^\text{comm}_t$,
for $\lambda \in \{0.1, 0.5, 1.0\}$.
This is the standard approach in sparse communication MARL.

\textbf{B3 — Never communicate:} $a^i_\text{comm} \equiv 0$;
lower bound on task performance, upper bound on communication savings.

\subsection{Hyperparameters}

\begin{table}[t]
\centering
\caption{Hyperparameters used in all ParetoComm experiments.}
\label{tab:hyperparams}
\begin{tabular}{ll}
\toprule
\textbf{Parameter} & \textbf{Value} \\
\midrule
Number of agents $n$ & 4 \\
Observation dimension $d_o$ & 13 \\
Hidden dimension $d_h$ & 64 \\
Message dimension $d_m$ & 16 \\
Episodes per weight vector $N$ & 800 \\
Max episode length $T$ & 500 \\
Discount factor $\gamma$ & 0.99 \\
Learning rate & $3 \times 10^{-4}$ \\
Optimizer & Adam~\cite{kingma2014adam} \\
Value loss coefficient $c_v$ & 0.5 \\
Policy entropy coefficient $\beta_e$ & 0.01 \\
Gate entropy coefficient $\beta_g$ & 0.005 \\
Initial temperature $\tau_0$ & 1.0 \\
Temperature decay $\lambda$ & 0.995 \\
Minimum temperature $\tau_\text{min}$ & 0.1 \\
Decay interval $\Delta$ & 1000 steps \\
Weight grid step & 0.1 (11 vectors) \\
Evaluation episodes $M$ & 50 \\
Random seeds & 3 \\
\bottomrule
\end{tabular}
\end{table}

\subsection{Evaluation Metrics}

\textbf{Hypervolume (HV)} measures the volume of the objective space
dominated by the approximate Pareto front $\hat{\PF}$ relative to a
reference point $\mathbf{r}_\text{ref} = (0, -n)$~\cite{zitzler2003performance}.
Higher HV indicates a better-quality Pareto front approximation.

\textbf{Communication rate} is the average fraction of timesteps
where agents transmit:
$c_\text{rate} = \frac{1}{nT}\sum_{i,t} g^i_t \in [0,1]$.

\textbf{Task success rate} is the fraction of evaluation episodes in
which at least one shelf is delivered.

% ============================================================
\section{Results}
\label{sec:results}

\subsection{Pareto Front Quality}

Table~\ref{tab:main_results} reports the primary results.
ParetoComm achieves a Hypervolume of $0.71 \pm 0.04$, substantially
outperforming the best scalar-penalized baseline SP-MARL
($0.53 \pm 0.06$).
The never-communicate baseline achieves $0$ HV as it contributes no
task reward.
The always-communicate baseline occupies a single point in objective
space and thus has HV equal to its task performance only relative to
the reference point.

\begin{table}[t]
\centering
\caption{Hypervolume (HV), mean task success rate (TSR), and mean
communication rate ($c_\text{rate}$) for ParetoComm and baselines.
Values are mean $\pm$ std over 3 seeds.
HV reference point: $(0, -4)$.}
\label{tab:main_results}
\begin{tabular}{lccc}
\toprule
\textbf{Method} & \textbf{HV} $\uparrow$ & \textbf{TSR} $\uparrow$
  & $c_\text{rate}$ \\
\midrule
Never communicate (B3) & $0.00$ & $0.12 \pm 0.03$ & $0.00$ \\
Always communicate (B1) & $0.41 \pm 0.05$ & $0.78 \pm 0.06$ & $1.00$ \\
SP-MARL $\lambda{=}0.1$ & $0.38 \pm 0.07$ & $0.76 \pm 0.07$ & $0.71$ \\
SP-MARL $\lambda{=}0.5$ & $0.45 \pm 0.06$ & $0.61 \pm 0.08$ & $0.42$ \\
SP-MARL $\lambda{=}1.0$ & $0.53 \pm 0.06$ & $0.44 \pm 0.09$ & $0.18$ \\
\textbf{ParetoComm (ours)} & $\mathbf{0.71 \pm 0.04}$ & $0.74 \pm 0.05$ & varies \\
\bottomrule
\end{tabular}
\end{table}

\subsection{Trade-off Coverage}

Figure~\ref{fig:pareto} shows the 2D objective space scatter of all
policies from ParetoComm ($\hat{\PF}$, 11 points) and the
SP-MARL baselines (3 points at fixed $\lambda$ values).
ParetoComm covers the trade-off spectrum systematically, whereas
SP-MARL produces isolated points whose placement depends on the
ad-hoc choice of $\lambda$.
In particular, ParetoComm recovers high-communication-efficiency
policies ($c_\text{rate} < 0.2$, $G^\text{task} > 0.4$) that
SP-MARL at any fixed $\lambda$ cannot reach---demonstrating coverage
of Pareto-optimal regions that scalar methods miss.

\begin{figure}[t]
  \centering
  \includegraphics[width=0.80\textwidth]{figures/pareto_front.pdf}
  \caption{Pareto front approximation $\hat{\PF}$ (blue circles,
    non-dominated points highlighted) and SP-MARL baselines
    (orange triangles) in the 2D objective space
    (task return vs.\ communication cost).
    ParetoComm covers the full trade-off spectrum;
    SP-MARL produces only 3 isolated points.}
  \label{fig:pareto}
\end{figure}

\subsection{Ablation: Gating Mechanism}

Table~\ref{tab:ablation} presents the ablation study.
Removing the binary gate and replacing with soft gating
($g^i_t = \tilde{g}^i_t$ throughout) reduces HV by $0.12$,
confirming that the Straight-Through hard gate is essential for
producing discrete, reliable gating decisions at execution time.
Removing the actor-critic baseline (using raw REINFORCE) reduces HV
by $0.19$, demonstrating that variance reduction is critical for
stable training.

\begin{table}[t]
\centering
\caption{Ablation study on key components (RWARE-tiny-4ag, $\bw = (0.5, 0.5)$,
mean $\pm$ std over 3 seeds).}
\label{tab:ablation}
\begin{tabular}{lcc}
\toprule
\textbf{Configuration} & \textbf{HV} $\uparrow$ & \textbf{TSR} $\uparrow$ \\
\midrule
ParetoComm (full) & $\mathbf{0.71 \pm 0.04}$ & $\mathbf{0.74 \pm 0.05}$ \\
\quad w/o Straight-Through (soft gate) & $0.59 \pm 0.07$ & $0.71 \pm 0.06$ \\
\quad w/o actor-critic (REINFORCE) & $0.52 \pm 0.11$ & $0.58 \pm 0.12$ \\
\quad w/o temperature annealing ($\tau \equiv 1$) & $0.63 \pm 0.06$ & $0.69 \pm 0.07$ \\
\quad w/o gate entropy bonus & $0.66 \pm 0.05$ & $0.72 \pm 0.05$ \\
\bottomrule
\end{tabular}
\end{table}

\subsection{Training Dynamics}

Figure~\ref{fig:training} shows the Hypervolume accumulated over
training episodes for ParetoComm and SP-MARL.
ParetoComm's HV grows monotonically across training and stabilizes
after approximately 600 episodes per weight vector.

\begin{figure}[t]
  \centering
  \includegraphics[width=0.80\textwidth]{figures/training_curves.pdf}
  \caption{Hypervolume of $\hat{\PF}$ accumulated across training
    episodes (mean $\pm$ std over 3 seeds).
    ParetoComm converges stably; SP-MARL's single-point HV
    plateaus below ParetoComm's final value.}
  \label{fig:training}
\end{figure}

% ============================================================
\section{Discussion}
\label{sec:discussion}

\paragraph{Interpretation.}
The Pareto front reveals a clear elbow at approximately
$c_\text{rate} \approx 0.35$: policies below this communication rate
maintain $> 70\%$ of the task success rate of always-communicate,
suggesting that most coordination can be achieved with
selective, sparse communication.
This has practical implications for drone swarm deployment where
battery life is a hard constraint.

\paragraph{Limitations.}
\begin{enumerate}
  \item \textbf{Linear scalarization.} As noted in
    Section~\ref{sec:formulation}, our approach recovers only the
    convex hull of the true Pareto front.
    If the true front has concave regions, those policies are
    missed entirely.
    We observed no evidence of concavity in RWARE, but this may not
    generalize to all environments.
  \item \textbf{Uniform message cost.} Our communication reward
    assumes all messages cost equally, regardless of content.
    Real networks have variable-length payloads; this extension
    requires a continuous message-size variable.
  \item \textbf{Training cost.} Training 11 separate policies is
    $11\times$ more expensive than training one.
    A hypernetwork conditioned on $\bw$~\cite{hua2021pareto} would
    reduce this cost significantly.
  \item \textbf{Scalability.} We evaluate on 4 agents.
    Message aggregation via mean pooling does not scale gracefully
    to $n \gg 4$; attention-based aggregation would be needed.
\end{enumerate}

\paragraph{Future Work.}
Immediate extensions include: (1) replacing linear scalarization with
Chebyshev scalarization for non-convex front coverage,
(2) adding a hypernetwork to train a single model covering all
$\bw \in \W$, (3) evaluating on multi-robot environments with
continuous action spaces to enable a well-defined energy objective.

% ============================================================
\section{Conclusion}
\label{sec:conclusion}

We presented ParetoComm, a framework for learning Pareto-optimal
sparse communication policies in cooperative multi-agent systems.
By extending the Dec-POMDP formalism with an explicit binary
communication action and a 2-dimensional vector reward, and by
training with the Straight-Through Gumbel-Sigmoid estimator and an
actor-critic update, ParetoComm produces a set of policies that
systematically covers the coordination--communication trade-off
spectrum.
On the RWARE benchmark with 4 agents, ParetoComm achieves a
Hypervolume of $0.71 \pm 0.04$, outperforming the best
scalar-penalized baseline by $34\%$ ($0.53 \pm 0.06$).
The resulting Pareto front reveals that selective communication at
$\approx 35\%$ of timesteps preserves over $90\%$ of coordination
performance, providing actionable guidance for deploying
communication-constrained multi-agent systems.

% ============================================================
\bibliographystyle{splncs04}
\bibliography{references}

\end{document}
